\scanpage{68}
\mspage{- 9 -}
\margin{n° 384}
Puisque $Q_\alpha$ est compact, la constante 2 appartient à $\mathcal K_+(Q_\alpha)$ et il existe une fonction $g\in\mathcal K_+(G)$ avec $g*\alpha(x)=2$ pour tout $x\in G$. Soit $K$ le support de $g$ et posons $L=DK$. Il existe un voisinage $W$ de $\alpha$ dans $T^\circ$ tel que $g*\beta(x)\geq 1$ pour $x\in L$ et $\beta\in W$. Soit $\beta\in W$. Comme le support de $g*\beta$ est contenu dans $KH_\beta$, on a donc $DK\subset KH_\beta$. Supposons démontré que $D^{r-1}K\subset KH_\beta$ : on a alors :
\[
(14)\qquad D^rK=DD^{r-1}K\subset DKH_\beta\subset KH_\beta H_\beta=KH_\beta.
\]
Comme $G$ est la réunion des $D^r$, on en déduit que $G=KH_\beta$. De plus, on a $g*\beta\geq 1$ sur $K$ et $g*\beta$ est invariante à droite par $H_\beta$ : on a donc $g*\beta\geq 1$ partout et ceci montre que $W\subset T^\circ(g)$. Donc $T^\circ(g)$ est bien un voisinage de $\alpha$ et il est clair que $\beta\in T^\circ(g)$ entraîne $G=KH_\beta$ et par suite que $G/H_\beta$ est compact.

COROLLAIRE 1.- \uline{L'ensemble} $T^c$ \uline{des} $\alpha\in T^\circ$ \uline{telles que} $G/H_\alpha$ \uline{soit compact est un ouvert de} $T^\circ$ \uline{et la restriction de} $v$ \uline{à} $T^c$ \uline{est continue}.

COROLLAIRE 2.- \uline{L'ensemble} $S^c$ \uline{des sous-groupes fermés unimodulaires à quotient compact est un sous-ensemble ouvert de l'ensemble compact} $S^\circ$ \uline{des sous-groupes} \inserted{\uline{fermés unimodulaires}.}

La proposition 6 n'est pas exacte si on ne suppose pas que $G$ est engendré par un compact (exer.\ ).

7. \uline{Sous-groupes discrets.}

Soit $D$ le sous-espace de $S$ formé des sous-groupes \uline{discrets} de $G$ et soit $\Delta$ le sous-espace de $T$ formé des mesures de Haar \uline{normalisées} des éléments de $D$ : pour qu'un élément $\alpha\in T$ appartienne à $\Delta$, il faut et il suffit que $\alpha(\{e\})=1$. On a évidemment $D\subset S^\circ$ et $\Delta\subset T^\circ$.

Pour tout compact $K$ de $G$, l'application $\alpha\longrightarrow\alpha(K)$ est semi-continue \uline{supérieurement} sur $\mathfrak M_+(G)$ (ch. IV, § 4, n° 4, Remarque). Par suite l'ensemble des $\alpha\in T$ telles que $\alpha(\{e\})\geq 1$ est \uline{fermé} dans $T$. D'autre part, pour tout voisinage ouvert $U$ de $e$, l'application $\alpha\longrightarrow\alpha(U)$ est semi-continue inférieurement sur
\margin{67}
